\bookchapter{第五十二章：四大組件同 Activity 生命週期}

\chapterintro{如果話 Compose 係 App 嘅「面」，咁呢一章講嘅就係 App 嘅「骨架」。安卓系統有四大基本組件，所有 App 都係由佢哋砌出嚟嘅；而當中最重要嘅 Activity，有一張每個安卓開發者都必須刻入腦嘅「生命週期圖」。本章學完，你就可以理解「點解切咗去後台返嚟，我個 App 嘅狀態就冇咗」呢類千古之謎。}

\booksection{一、四大組件：安卓嘅「四樑八柱」}

安卓官方定義咗四種基本組件，分別負責一類工作：

\begin{center}
\begin{tabularx}{\linewidth}{|X|X|X|}
\hline
\textbf{組件} & \textbf{職責} & \textbf{打個比方} \tabularnewline
\hline
Activity & 一個睇得到嘅介面 & 舖頭嘅門面——用戶見到、又互動到 \tabularnewline
Service & 喺後台做嘢嘅「隱形人」 & 倉庫工人——播音樂、下載檔案，唔需要介面 \tabularnewline
BroadcastReceiver & 接收系統廣播 & 傳達室——「冇電喇」「斷網喇」嘅喇叭 \tabularnewline
ContentProvider & 跨 App 共享資料 & 公共資料櫃——例如將聯絡人共享畀其他 App \tabularnewline
\hline
\end{tabularx}
\end{center}

\begin{center}\small\color{gray}安卓四大組件速覽\end{center}

新手階段你 95\% 嘅時間都係同 Activity 打交道；Service 同 BroadcastReceiver 做音樂播放、訊息推送嗰陣會用到；ContentProvider 用得最少（系統已經將常用資料封裝好）。另外留意：四大組件全部都要喺 AndroidManifest.xml 入面「登記戶口」，唔登記嘅話系統就搵唔到佢哋。

\booksection{二、Activity 生命週期：安卓開發嘅第一張「藏寶圖」}

手機上面嘅 App 隨時可能被打斷：嚟電話喇、用戶切咗去覆微信、系統記憶體唔夠將你殺咗……Activity 嘅一生，就係喺呢張圖度循環旅行：

onCreate（出生：初始化）

onStart（可見了）

onResume（前臺交互中）

onPause（部分被遮住）

onStop（完全看不見）

onDestroy（銷燬）

onRestart / onStart / onResume：回到前臺時原路返回

\begin{center}\small\color{gray}圖 4：Activity 生命週期（前進方向：上排出世 下排離場，返嚟嗰陣原路折返）\end{center}

\begin{center}
\begin{tabularx}{\linewidth}{|X|X|X|}
\hline
\textbf{回調} & \textbf{幾時觸發} & \textbf{典型用途} \tabularnewline
\hline
onCreate & Activity 開嗰陣（一次） & 載入介面、初始化資料 \tabularnewline
onStart & 介面變為可見 & 開始啲輕量邏輯 \tabularnewline
onResume & 攞到焦點、可以互動 & 開始動畫、恢復播放 \tabularnewline
onPause & 被部分遮住／切走嗰一刻 & 暫停動畫、儲存草稿 \tabularnewline
onStop & 完全睇唔到 & 釋放重磅資源、暫停下載 \tabularnewline
onDestroy & 被銷毀之前 & 最後嘅清理 \tabularnewline
\hline
\end{tabularx}
\end{center}

\begin{center}\small\color{gray}六個生命週期回調嘅典型用途\end{center}

\textbf{避坑提醒：轉螢幕會「殺死」Activity}\par

新手最經典嘅炒車現場：App 轉一轉螢幕，輸入框入面嘅字全部冇晒。原因：轉螢幕預設會銷毀並重建 Activity（行一次 onDestroy onCreate）。解法：用 rememberSaveable 儲存狀態（Compose 入面代替 remember），佢會喺重建嗰陣自動將個值搵返。第 6、7 章會反覆用到。

\booksection{三、Intent：App 內外通行嘅「車飛」}

Intent（意圖）係組件之間通訊嘅「車飛」。你想由一個頁面跳去另一個頁面，或者打開瀏覽器、叫起相機，都要「買飛上車」：

\begin{lstlisting}[language=Kotlin]
// 方式一：顯式 Intent——明確指定去邊個頁面
val intent = Intent(this, DetailActivity::class.java)
intent.putExtra("name", "張三")     // 車上帶埋"行李"（數據）
startActivity(intent)

// 方式二：隱式 Intent——淨係講要做乜，系統搵邊個做
val web = Intent(Intent.ACTION_VIEW, Uri.parse("https://developer.android.com"))
startActivity(web)   // 系統會開瀏覽器／或者畀你揀
\end{lstlisting}

兩種方式各有用途：顯式用於自己 App 入面嘅頁面跳轉；隱式用於調用系統或者其他 App 嘅能力（開網頁、打電話、分享）。Compose 項目入面，喺 Composable 入面用 LocalContext.current 攞到 Context 就可以 startActivity。

\booksection{四、多頁面同導航}

真實嘅 App 都有好多個頁面。Compose 官方推薦嘅導航方案係 Navigation Compose（依賴庫 androidx.navigation）：

\begin{lstlisting}[language=Kotlin]
val navController = rememberNavController()

NavHost(navController, startDestination = "home") {
    composable("home") { HomeScreen(navController) }
    composable("detail/{name}") { entry ->
        val name = entry.arguments?.getString("name")
        DetailScreen(name ?: "未知")
    }
}

// 喺頁面入面跳轉：
navController.navigate("detail/張三")     // 帶參數
navController.popBackStack()              // 返回上一頁
\end{lstlisting}

\begin{tcbinfo}{小貼士：頁面多咗點算？}
項目大咗之後，將所有頁面路由集中寫喺一個檔案入面（例如 AppNavGraph.kt），統一管理，唔好畀「跳去邊」嘅邏輯散落喺各個頁面。呢一步係由「行到」走向「可維護」嘅第一步。

\end{tcbinfo}

\booksection{五、順便認識：Fragment 係乜嘢（同埋點解本書冇教）}

查舊教程嗰陣你一定會撞到 Fragment——佢係 Activity 入面嘅「可重用介面片段」，喺舊式多頁面架構入面係絕對主角。但係 Compose 年代，頁面直接用 Composable 函數 + Navigation 路由就夠，官方亦都已經淡化咗 Fragment 嘅地位。本書因此冇將佢放入主線。

畀你嘅建議：新項目唔好主動用 Fragment；但係睇舊程式碼、面試被問到嗰陣，知道佢係「Activity 入面嵌住嘅介面片段」就得。學習資源永遠優先揀 2023 年之後嘅——安卓技術棧嘅「老黃曆」更新得比想像中快。

\begin{tcbpractice}{嚟動手：畀計數器裝上「記憶」}
將第 51 章嘅計數器改一改：(1) 將 remember 換成 rememberSaveable，轉螢幕驗證個數字仲喺唔喺度；(2) 加一個「關於」按鈕，用 Navigation 跳去第二個頁面顯示版本號；(3) 喺兩個頁面嘅 onCreate/onPause 入面 print 日誌（Log.d），用 Logcat 觀察生命週期調用次序。做完呢三步，生命週期就唔再係紙上談兵。

\end{tcbpractice}

\begin{tcbsummary}{本章小結}
\begin{itemize}
  \item 四大組件：Activity（門面）、Service（後台）、BroadcastReceiver（喇叭）、ContentProvider（資料櫃）
  \item 生命週期六回調：onCreate onStart onResume onPause onStop onDestroy
  \item 轉螢幕會銷毀重建 Activity：用 rememberSaveable 保住狀態
  \item Intent 係「車飛」：顯式跳自己頁面，隱式調系統能力
  \item 多頁面用 Navigation Compose：NavHost + composable 路由
\end{itemize}

\end{tcbsummary}
